\documentclass[10pt,aspectratio=169]{beamer}
\usetheme{Madrid}
\usecolortheme{default}
\usepackage[utf8]{inputenc}
\usepackage[T1]{fontenc}
\usepackage[spanish,mexico]{babel}
\usepackage{amsmath,amssymb,mathtools}
\usepackage{tikz}
\usepackage{booktabs,array}
\setbeamertemplate{navigation symbols}{}
\setbeamertemplate{footline}[frame number]
\setbeamertemplate{theorems}[numbered]
\newtheorem{observacion}{Observación}
\newtheorem{pregunta}{Pregunta guía}

\newcommand{\ejercicio}[1]{%
  \begin{center}
  \vfill
  {\fontsize{24}{30}\selectfont #1}
  \vfill
  \end{center}
}

\title{Cálculo Integral}
\subtitle{Integrales indefinidas: ejercicios para trabajo en clase}
\author{Carlos Ernesto Martínez}
\institute{Tecnológico Nacional de México}
\date{}
\begin{document}
\begin{frame}\titlepage\end{frame}
\section{Integración inmediata: regla de la potencia}
\begin{frame}{Integración inmediata: regla de la potencia}
\small Ejercicio 1
\ejercicio{$\displaystyle \int x^7\,dx$}
\end{frame}

\begin{frame}{Integración inmediata: regla de la potencia}
\small Ejercicio 2
\ejercicio{$\displaystyle \int x^{12}\,dx$}
\end{frame}

\begin{frame}{Integración inmediata: regla de la potencia}
\small Ejercicio 3
\ejercicio{$\displaystyle \int x^{-3}\,dx$}
\end{frame}

\begin{frame}{Integración inmediata: regla de la potencia}
\small Ejercicio 4
\ejercicio{$\displaystyle \int x^{-5}\,dx$}
\end{frame}

\begin{frame}{Integración inmediata: regla de la potencia}
\small Ejercicio 5
\ejercicio{$\displaystyle \int x^{1/2}\,dx$}
\end{frame}

\begin{frame}{Integración inmediata: regla de la potencia}
\small Ejercicio 6
\ejercicio{$\displaystyle \int x^{3/2}\,dx$}
\end{frame}

\begin{frame}{Integración inmediata: regla de la potencia}
\small Ejercicio 7
\ejercicio{$\displaystyle \int x^{-1/2}\,dx$}
\end{frame}

\begin{frame}{Integración inmediata: regla de la potencia}
\small Ejercicio 8
\ejercicio{$\displaystyle \int x^{5/3}\,dx$}
\end{frame}

\section{Polinomios y combinaciones lineales}
\begin{frame}{Polinomios y combinaciones lineales}
\small Ejercicio 1
\ejercicio{$\displaystyle \int(3x^2-4x+7)\,dx$}
\end{frame}

\begin{frame}{Polinomios y combinaciones lineales}
\small Ejercicio 2
\ejercicio{$\displaystyle \int(5x^4-2x^3+x-6)\,dx$}
\end{frame}

\begin{frame}{Polinomios y combinaciones lineales}
\small Ejercicio 3
\ejercicio{$\displaystyle \int(2x^5+7x^2-9)\,dx$}
\end{frame}

\begin{frame}{Polinomios y combinaciones lineales}
\small Ejercicio 4
\ejercicio{$\displaystyle \int(8-3x+4x^3)\,dx$}
\end{frame}

\begin{frame}{Polinomios y combinaciones lineales}
\small Ejercicio 5
\ejercicio{$\displaystyle \int\left(\frac{x^4}{2}-\frac{3x^2}{4}+5x\right)dx$}
\end{frame}

\begin{frame}{Polinomios y combinaciones lineales}
\small Ejercicio 6
\ejercicio{$\displaystyle \int(6x^7-5x^5+2x^2-1)\,dx$}
\end{frame}

\begin{frame}{Polinomios y combinaciones lineales}
\small Ejercicio 7
\ejercicio{$\displaystyle \int\left(\frac{2}{3}x^6-\frac{5}{2}x^3+4\right)dx$}
\end{frame}

\begin{frame}{Polinomios y combinaciones lineales}
\small Ejercicio 8
\ejercicio{$\displaystyle \int(7x^8-3x^4+x^2-8x+11)\,dx$}
\end{frame}

\section{Radicales y potencias fraccionarias}
\begin{frame}{Radicales y potencias fraccionarias}
\small Ejercicio 1
\ejercicio{$\displaystyle \int \sqrt{x}\,dx$}
\end{frame}

\begin{frame}{Radicales y potencias fraccionarias}
\small Ejercicio 2
\ejercicio{$\displaystyle \int \sqrt[3]{x^2}\,dx$}
\end{frame}

\begin{frame}{Radicales y potencias fraccionarias}
\small Ejercicio 3
\ejercicio{$\displaystyle \int \frac{1}{\sqrt{x}}\,dx$}
\end{frame}

\begin{frame}{Radicales y potencias fraccionarias}
\small Ejercicio 4
\ejercicio{$\displaystyle \int \frac{1}{\sqrt[3]{x^2}}\,dx$}
\end{frame}

\begin{frame}{Radicales y potencias fraccionarias}
\small Ejercicio 5
\ejercicio{$\displaystyle \int(2\sqrt{x}+3\sqrt[3]{x})\,dx$}
\end{frame}

\begin{frame}{Radicales y potencias fraccionarias}
\small Ejercicio 6
\ejercicio{$\displaystyle \int\frac{x^2+\sqrt{x}}{x}\,dx$}
\end{frame}

\begin{frame}{Radicales y potencias fraccionarias}
\small Ejercicio 7
\ejercicio{$\displaystyle \int\left(\sqrt[4]{x^3}+\frac{2}{\sqrt{x}}\right)dx$}
\end{frame}

\begin{frame}{Radicales y potencias fraccionarias}
\small Ejercicio 8
\ejercicio{$\displaystyle \int\frac{x^3+2x}{\sqrt{x}}\,dx$}
\end{frame}

\section{Simplificación algebraica antes de integrar}
\begin{frame}{Simplificación algebraica antes de integrar}
\small Ejercicio 1
\ejercicio{$\displaystyle \int\frac{x^3+2x^2-x}{x}\,dx$}
\end{frame}

\begin{frame}{Simplificación algebraica antes de integrar}
\small Ejercicio 2
\ejercicio{$\displaystyle \int\frac{3x^4-2x^2+5x}{x^2}\,dx$}
\end{frame}

\begin{frame}{Simplificación algebraica antes de integrar}
\small Ejercicio 3
\ejercicio{$\displaystyle \int\frac{(x+1)^2}{x}\,dx$}
\end{frame}

\begin{frame}{Simplificación algebraica antes de integrar}
\small Ejercicio 4
\ejercicio{$\displaystyle \int\frac{x^4-1}{x^2}\,dx$}
\end{frame}

\begin{frame}{Simplificación algebraica antes de integrar}
\small Ejercicio 5
\ejercicio{$\displaystyle \int\frac{2x^3+5x^2-3x}{x}\,dx$}
\end{frame}

\begin{frame}{Simplificación algebraica antes de integrar}
\small Ejercicio 6
\ejercicio{$\displaystyle \int\frac{x^5+x^3+x}{x^2}\,dx$}
\end{frame}

\begin{frame}{Simplificación algebraica antes de integrar}
\small Ejercicio 7
\ejercicio{$\displaystyle \int\frac{(x-2)(x+3)}{x}\,dx$}
\end{frame}

\begin{frame}{Simplificación algebraica antes de integrar}
\small Ejercicio 8
\ejercicio{$\displaystyle \int\frac{x^6-2x^3+x}{x^2}\,dx$}
\end{frame}

\section{La forma $1/x$ y logaritmos inmediatos}
\begin{frame}{La forma $1/x$ y logaritmos inmediatos}
\small Ejercicio 1
\ejercicio{$\displaystyle \int\frac{1}{x}\,dx$}
\end{frame}

\begin{frame}{La forma $1/x$ y logaritmos inmediatos}
\small Ejercicio 2
\ejercicio{$\displaystyle \int\frac{5}{x}\,dx$}
\end{frame}

\begin{frame}{La forma $1/x$ y logaritmos inmediatos}
\small Ejercicio 3
\ejercicio{$\displaystyle \int\left(3x^2+\frac{2}{x}\right)dx$}
\end{frame}

\begin{frame}{La forma $1/x$ y logaritmos inmediatos}
\small Ejercicio 4
\ejercicio{$\displaystyle \int\left(\frac{4}{x}-x^3\right)dx$}
\end{frame}

\begin{frame}{La forma $1/x$ y logaritmos inmediatos}
\small Ejercicio 5
\ejercicio{$\displaystyle \int\frac{7x^2+3}{x}\,dx$}
\end{frame}

\begin{frame}{La forma $1/x$ y logaritmos inmediatos}
\small Ejercicio 6
\ejercicio{$\displaystyle \int\frac{x^3-4}{x}\,dx$}
\end{frame}

\begin{frame}{La forma $1/x$ y logaritmos inmediatos}
\small Ejercicio 7
\ejercicio{$\displaystyle \int\left(\sqrt{x}+\frac{6}{x}\right)dx$}
\end{frame}

\begin{frame}{La forma $1/x$ y logaritmos inmediatos}
\small Ejercicio 8
\ejercicio{$\displaystyle \int\frac{2x^4-x+5}{x}\,dx$}
\end{frame}

\section{Funciones exponenciales}

\begin{frame}{Funciones exponenciales: dos casos fundamentales}
\small
Para la base natural,
\[
\frac{d}{dx}e^x=e^x
\qquad\Longrightarrow\qquad
\boxed{\int e^x\,dx=e^x+C}.
\]
\pause
Para una base general $a>0$, $a\neq1$,
\[
a^x=e^{x\ln a}.
\]
Por ello, la derivada contiene un factor adicional:
\[
\boxed{\frac{d}{dx}a^x=a^x\ln a}.
\]
\end{frame}

\begin{frame}{Deducción de $\int a^x\,dx$}
\small
Partimos de
\[
\int a^x\,dx=\int e^{x\ln a}\,dx.
\]
Hacemos
\[
u=x\ln a,
\qquad du=\ln a\,dx,
\qquad dx=\frac{du}{\ln a}.
\]
Entonces
\[
\begin{aligned}
\int a^x\,dx
&=\frac{1}{\ln a}\int e^u\,du\\
&=\frac{e^u}{\ln a}+C\\
&=\boxed{\frac{a^x}{\ln a}+C}.
\end{aligned}
\]
\end{frame}

\begin{frame}{Ejemplo guiado: $\int 2^x\,dx$}
\small
Aplicando la fórmula general,
\[
\boxed{\int 2^x\,dx=\frac{2^x}{\ln2}+C}.
\]
\pause
Verificación:
\[
\frac{d}{dx}\left(\frac{2^x}{\ln2}\right)
=\frac{1}{\ln2}\,2^x\ln2
=2^x.
\]
\pause
\begin{block}{Idea clave}
No se debe escribir $\int 2^x\,dx=2^x+C$: el factor $\ln2$ debe compensarse.
\end{block}
\end{frame}

\begin{frame}{Exponente lineal: $a^{mx+b}$}
\small
Sea
\[
\int a^{mx+b}\,dx.
\]
Con
\[
u=mx+b,\qquad du=m\,dx,
\]
se obtiene
\[
\boxed{\int a^{mx+b}\,dx
=\frac{a^{mx+b}}{m\ln a}+C}.
\]
\pause
Por ejemplo,
\[
\int2^{3x+1}\,dx
=\boxed{\frac{2^{3x+1}}{3\ln2}+C}.
\]
\end{frame}

\begin{frame}{Exponenciales y sustitución}
\small
Si el exponente es una función $g(x)$ y aparece su derivada,
\[
\boxed{\int a^{g(x)}g'(x)\,dx
=\frac{a^{g(x)}}{\ln a}+C}.
\]
\pause
Ejemplo:
\[
\int x\,2^{x^2}\,dx.
\]
Tomando $u=x^2$, $du=2x\,dx$,
\[
\int x\,2^{x^2}\,dx
=\frac12\int2^u\,du
=\boxed{\frac{2^{x^2}}{2\ln2}+C}.
\]
\end{frame}

\begin{frame}{Funciones exponenciales}
\small Ejercicio 1
\ejercicio{$\displaystyle \int e^x\,dx$}
\end{frame}

\begin{frame}{Funciones exponenciales}
\small Ejercicio 2
\ejercicio{$\displaystyle \int e^{3x}\,dx$}
\end{frame}

\begin{frame}{Funciones exponenciales}
\small Ejercicio 3
\ejercicio{$\displaystyle \int e^{-2x}\,dx$}
\end{frame}

\begin{frame}{Funciones exponenciales}
\small Ejercicio 4
\ejercicio{$\displaystyle \int 4e^{5x}\,dx$}
\end{frame}

\begin{frame}{Funciones exponenciales}
\small Ejercicio 5
\ejercicio{$\displaystyle \int 2^x\,dx$}
\end{frame}

\begin{frame}{Funciones exponenciales}
\small Ejercicio 6
\ejercicio{$\displaystyle \int 5^{3x}\,dx$}
\end{frame}

\begin{frame}{Funciones exponenciales}
\small Ejercicio 7
\ejercicio{$\displaystyle \int 3^{-2x}\,dx$}
\end{frame}

\begin{frame}{Funciones exponenciales}
\small Ejercicio 8
\ejercicio{$\displaystyle \int(2e^{4x}-3\,7^x)\,dx$}
\end{frame}

\begin{frame}{Funciones exponenciales}
\small Ejercicio 9
\ejercicio{$\displaystyle \int 10^x\,dx$}
\end{frame}

\begin{frame}{Funciones exponenciales}
\small Ejercicio 10
\ejercicio{$\displaystyle \int 5^{2x-3}\,dx$}
\end{frame}

\begin{frame}{Funciones exponenciales}
\small Ejercicio 11
\ejercicio{$\displaystyle \int x\,4^{x^2}\,dx$}
\end{frame}

\begin{frame}{Funciones exponenciales}
\small Ejercicio 12
\ejercicio{$\displaystyle \int (2x+1)\,5^{x^2+x}\,dx$}
\end{frame}

\begin{frame}{Funciones exponenciales}
\small Ejercicio 13
\ejercicio{$\displaystyle \int \cos x\,2^{\sin x}\,dx$}
\end{frame}

\section{Trigonométricas inmediatas}
\begin{frame}{Trigonométricas inmediatas}
\small Ejercicio 1
\ejercicio{$\displaystyle \int\sin x\,dx$}
\end{frame}

\begin{frame}{Trigonométricas inmediatas}
\small Ejercicio 2
\ejercicio{$\displaystyle \int\cos x\,dx$}
\end{frame}

\begin{frame}{Trigonométricas inmediatas}
\small Ejercicio 3
\ejercicio{$\displaystyle \int\sec^2x\,dx$}
\end{frame}

\begin{frame}{Trigonométricas inmediatas}
\small Ejercicio 4
\ejercicio{$\displaystyle \int\csc^2x\,dx$}
\end{frame}

\begin{frame}{Trigonométricas inmediatas}
\small Ejercicio 5
\ejercicio{$\displaystyle \int\sec x\tan x\,dx$}
\end{frame}

\begin{frame}{Trigonométricas inmediatas}
\small Ejercicio 6
\ejercicio{$\displaystyle \int\csc x\cot x\,dx$}
\end{frame}

\begin{frame}{Trigonométricas inmediatas}
\small Ejercicio 7
\ejercicio{$\displaystyle \int(3\cos x-2\sin x)\,dx$}
\end{frame}

\begin{frame}{Trigonométricas inmediatas}
\small Ejercicio 8
\ejercicio{$\displaystyle \int(4\sec^2x+5\csc x\cot x)\,dx$}
\end{frame}

\section{Trigonométricas con identidades elementales}
\begin{frame}{Trigonométricas con identidades elementales}
\small Ejercicio 1
\ejercicio{$\displaystyle \int\sin^2x\,dx$}
\end{frame}

\begin{frame}{Trigonométricas con identidades elementales}
\small Ejercicio 2
\ejercicio{$\displaystyle \int\cos^2x\,dx$}
\end{frame}

\begin{frame}{Trigonométricas con identidades elementales}
\small Ejercicio 3
\ejercicio{$\displaystyle \int(1+\cos 2x)\,dx$}
\end{frame}

\begin{frame}{Trigonométricas con identidades elementales}
\small Ejercicio 4
\ejercicio{$\displaystyle \int(1-\cos 2x)\,dx$}
\end{frame}

\begin{frame}{Trigonométricas con identidades elementales}
\small Ejercicio 5
\ejercicio{$\displaystyle \int\sin x\cos x\,dx$}
\end{frame}

\begin{frame}{Trigonométricas con identidades elementales}
\small Ejercicio 6
\ejercicio{$\displaystyle \int\sin(2x)\,dx$}
\end{frame}

\begin{frame}{Trigonométricas con identidades elementales}
\small Ejercicio 7
\ejercicio{$\displaystyle \int\cos(5x)\,dx$}
\end{frame}

\begin{frame}{Trigonométricas con identidades elementales}
\small Ejercicio 8
\ejercicio{$\displaystyle \int(3\sin^2x+2\cos^2x)\,dx$}
\end{frame}

\section{Sustitución: potencias de una función}
\begin{frame}{Sustitución: potencias de una función}
\small Ejercicio 1
\ejercicio{$\displaystyle \int 2x(x^2+1)^5\,dx$}
\end{frame}

\begin{frame}{Sustitución: potencias de una función}
\small Ejercicio 2
\ejercicio{$\displaystyle \int 3x^2(x^3-4)^6\,dx$}
\end{frame}

\begin{frame}{Sustitución: potencias de una función}
\small Ejercicio 3
\ejercicio{$\displaystyle \int(2x+1)(x^2+x+3)^4\,dx$}
\end{frame}

\begin{frame}{Sustitución: potencias de una función}
\small Ejercicio 4
\ejercicio{$\displaystyle \int 4x^3(x^4+2)^7\,dx$}
\end{frame}

\begin{frame}{Sustitución: potencias de una función}
\small Ejercicio 5
\ejercicio{$\displaystyle \int x(3x^2-1)^8\,dx$}
\end{frame}

\begin{frame}{Sustitución: potencias de una función}
\small Ejercicio 6
\ejercicio{$\displaystyle \int(6x-5)(3x^2-5x+2)^3\,dx$}
\end{frame}

\begin{frame}{Sustitución: potencias de una función}
\small Ejercicio 7
\ejercicio{$\displaystyle \int x^2(1-x^3)^5\,dx$}
\end{frame}

\begin{frame}{Sustitución: potencias de una función}
\small Ejercicio 8
\ejercicio{$\displaystyle \int(5x^4+2)(x^5+2x-1)^6\,dx$}
\end{frame}

\section{Sustitución con radicales}
\begin{frame}{Sustitución con radicales}
\small Ejercicio 1
\ejercicio{$\displaystyle \int\frac{x}{\sqrt{x^2+4}}\,dx$}
\end{frame}

\begin{frame}{Sustitución con radicales}
\small Ejercicio 2
\ejercicio{$\displaystyle \int\frac{2x+1}{\sqrt{x^2+x+5}}\,dx$}
\end{frame}

\begin{frame}{Sustitución con radicales}
\small Ejercicio 3
\ejercicio{$\displaystyle \int x^2\sqrt{x^3+1}\,dx$}
\end{frame}

\begin{frame}{Sustitución con radicales}
\small Ejercicio 4
\ejercicio{$\displaystyle \int\frac{x^3}{\sqrt{x^4+7}}\,dx$}
\end{frame}

\begin{frame}{Sustitución con radicales}
\small Ejercicio 5
\ejercicio{$\displaystyle \int(3x^2-2)\sqrt{x^3-2x+4}\,dx$}
\end{frame}

\begin{frame}{Sustitución con radicales}
\small Ejercicio 6
\ejercicio{$\displaystyle \int\frac{4x}{\sqrt[3]{2x^2+1}}\,dx$}
\end{frame}

\begin{frame}{Sustitución con radicales}
\small Ejercicio 7
\ejercicio{$\displaystyle \int x^2\sqrt[4]{x^3+2}\,dx$}
\end{frame}

\begin{frame}{Sustitución con radicales}
\small Ejercicio 8
\ejercicio{$\displaystyle \int\frac{5x^4}{\sqrt{1+x^5}}\,dx$}
\end{frame}

\section{Sustitución que conduce a logaritmos: $f'(x)/f(x)$}
\begin{frame}{Sustitución que conduce a logaritmos: $f'(x)/f(x)$}
\small Ejercicio 1
\ejercicio{$\displaystyle \int\frac{2x}{x^2+3}\,dx$}
\end{frame}

\begin{frame}{Sustitución que conduce a logaritmos: $f'(x)/f(x)$}
\small Ejercicio 2
\ejercicio{$\displaystyle \int\frac{3x^2}{x^3+5}\,dx$}
\end{frame}

\begin{frame}{Sustitución que conduce a logaritmos: $f'(x)/f(x)$}
\small Ejercicio 3
\ejercicio{$\displaystyle \int\frac{2x+1}{x^2+x+4}\,dx$}
\end{frame}

\begin{frame}{Sustitución que conduce a logaritmos: $f'(x)/f(x)$}
\small Ejercicio 4
\ejercicio{$\displaystyle \int\frac{4x^3}{x^4-7}\,dx$}
\end{frame}

\begin{frame}{Sustitución que conduce a logaritmos: $f'(x)/f(x)$}
\small Ejercicio 5
\ejercicio{$\displaystyle \int\frac{6x-2}{3x^2-2x+9}\,dx$}
\end{frame}

\begin{frame}{Sustitución que conduce a logaritmos: $f'(x)/f(x)$}
\small Ejercicio 6
\ejercicio{$\displaystyle \int\frac{5x^4+1}{x^5+x+2}\,dx$}
\end{frame}

\begin{frame}{Sustitución que conduce a logaritmos: $f'(x)/f(x)$}
\small Ejercicio 7
\ejercicio{$\displaystyle \int\frac{\cos x}{2+\sin x}\,dx$}
\end{frame}

\begin{frame}{Sustitución que conduce a logaritmos: $f'(x)/f(x)$}
\small Ejercicio 8
\ejercicio{$\displaystyle \int\frac{e^x}{1+e^x}\,dx$}
\end{frame}

\section{Sustitución con exponenciales}
\begin{frame}{Sustitución con exponenciales}
\small Ejercicio 1
\ejercicio{$\displaystyle \int e^x(1+e^x)^4\,dx$}
\end{frame}

\begin{frame}{Sustitución con exponenciales}
\small Ejercicio 2
\ejercicio{$\displaystyle \int e^{2x}\sqrt{1+e^{2x}}\,dx$}
\end{frame}

\begin{frame}{Sustitución con exponenciales}
\small Ejercicio 3
\ejercicio{$\displaystyle \int\frac{e^x}{(3+e^x)^2}\,dx$}
\end{frame}

\begin{frame}{Sustitución con exponenciales}
\small Ejercicio 4
\ejercicio{$\displaystyle \int e^{-x}(2+e^{-x})^5\,dx$}
\end{frame}

\begin{frame}{Sustitución con exponenciales}
\small Ejercicio 5
\ejercicio{$\displaystyle \int\frac{3e^{3x}}{1+e^{3x}}\,dx$}
\end{frame}

\begin{frame}{Sustitución con exponenciales}
\small Ejercicio 6
\ejercicio{$\displaystyle \int e^{4x}(1-e^{4x})^3\,dx$}
\end{frame}

\begin{frame}{Sustitución con exponenciales}
\small Ejercicio 7
\ejercicio{$\displaystyle \int\frac{2^x}{1+2^x}\,dx$}
\end{frame}

\begin{frame}{Sustitución con exponenciales}
\small Ejercicio 8
\ejercicio{$\displaystyle \int 5^x(2+5^x)^6\,dx$}
\end{frame}

\section{Sustitución con funciones trigonométricas}
\begin{frame}{Sustitución con funciones trigonométricas}
\small Ejercicio 1
\ejercicio{$\displaystyle \int\sin x\cos^4x\,dx$}
\end{frame}

\begin{frame}{Sustitución con funciones trigonométricas}
\small Ejercicio 2
\ejercicio{$\displaystyle \int\cos x\sin^5x\,dx$}
\end{frame}

\begin{frame}{Sustitución con funciones trigonométricas}
\small Ejercicio 3
\ejercicio{$\displaystyle \int\sec^2x(1+\tan x)^3\,dx$}
\end{frame}

\begin{frame}{Sustitución con funciones trigonométricas}
\small Ejercicio 4
\ejercicio{$\displaystyle \int\csc^2x(2-\cot x)^4\,dx$}
\end{frame}

\begin{frame}{Sustitución con funciones trigonométricas}
\small Ejercicio 5
\ejercicio{$\displaystyle \int\sec x\tan x(3+\sec x)^5\,dx$}
\end{frame}

\begin{frame}{Sustitución con funciones trigonométricas}
\small Ejercicio 6
\ejercicio{$\displaystyle \int\csc x\cot x(1+\csc x)^2\,dx$}
\end{frame}

\begin{frame}{Sustitución con funciones trigonométricas}
\small Ejercicio 7
\ejercicio{$\displaystyle \int\frac{\cos x}{1+\sin x}\,dx$}
\end{frame}

\begin{frame}{Sustitución con funciones trigonométricas}
\small Ejercicio 8
\ejercicio{$\displaystyle \int\frac{\sec^2x}{2+\tan x}\,dx$}
\end{frame}

\section{Integrales que conducen a funciones trigonométricas inversas}
\begin{frame}{Integrales que conducen a funciones trigonométricas inversas}
\small Ejercicio 1
\ejercicio{$\displaystyle \int\frac{1}{1+x^2}\,dx$}
\end{frame}

\begin{frame}{Integrales que conducen a funciones trigonométricas inversas}
\small Ejercicio 2
\ejercicio{$\displaystyle \int\frac{1}{\sqrt{1-x^2}}\,dx$}
\end{frame}

\begin{frame}{Integrales que conducen a funciones trigonométricas inversas}
\small Ejercicio 3
\ejercicio{$\displaystyle \int\frac{1}{4+x^2}\,dx$}
\end{frame}

\begin{frame}{Integrales que conducen a funciones trigonométricas inversas}
\small Ejercicio 4
\ejercicio{$\displaystyle \int\frac{1}{\sqrt{9-x^2}}\,dx$}
\end{frame}

\begin{frame}{Integrales que conducen a funciones trigonométricas inversas}
\small Ejercicio 5
\ejercicio{$\displaystyle \int\frac{1}{1+4x^2}\,dx$}
\end{frame}

\begin{frame}{Integrales que conducen a funciones trigonométricas inversas}
\small Ejercicio 6
\ejercicio{$\displaystyle \int\frac{1}{\sqrt{16-9x^2}}\,dx$}
\end{frame}

\begin{frame}{Integrales que conducen a funciones trigonométricas inversas}
\small Ejercicio 7
\ejercicio{$\displaystyle \int\frac{3}{1+9x^2}\,dx$}
\end{frame}

\begin{frame}{Integrales que conducen a funciones trigonométricas inversas}
\small Ejercicio 8
\ejercicio{$\displaystyle \int\frac{2}{\sqrt{25-4x^2}}\,dx$}
\end{frame}

\section{Integrales mixtas: identificar el método}
\begin{frame}{Integrales mixtas: identificar el método}
\small Ejercicio 1
\ejercicio{$\displaystyle \int(4x^3-2x+7)\,dx$}
\end{frame}

\begin{frame}{Integrales mixtas: identificar el método}
\small Ejercicio 2
\ejercicio{$\displaystyle \int\frac{x^2+3x}{x}\,dx$}
\end{frame}

\begin{frame}{Integrales mixtas: identificar el método}
\small Ejercicio 3
\ejercicio{$\displaystyle \int\frac{x}{x^2+9}\,dx$}
\end{frame}

\begin{frame}{Integrales mixtas: identificar el método}
\small Ejercicio 4
\ejercicio{$\displaystyle \int x^2(x^3+4)^7\,dx$}
\end{frame}

\begin{frame}{Integrales mixtas: identificar el método}
\small Ejercicio 5
\ejercicio{$\displaystyle \int\frac{x}{\sqrt{x^2+1}}\,dx$}
\end{frame}

\begin{frame}{Integrales mixtas: identificar el método}
\small Ejercicio 6
\ejercicio{$\displaystyle \int e^{2x}(1+e^{2x})^3\,dx$}
\end{frame}

\begin{frame}{Integrales mixtas: identificar el método}
\small Ejercicio 7
\ejercicio{$\displaystyle \int\cos x(2+\sin x)^5\,dx$}
\end{frame}

\begin{frame}{Integrales mixtas: identificar el método}
\small Ejercicio 8
\ejercicio{$\displaystyle \int\frac{\sec^2x}{1+\tan x}\,dx$}
\end{frame}

\begin{frame}{Integrales mixtas: identificar el método}
\small Ejercicio 9
\ejercicio{$\displaystyle \int\cos^2x\,dx$}
\end{frame}

\begin{frame}{Integrales mixtas: identificar el método}
\small Ejercicio 10
\ejercicio{$\displaystyle \int\frac{1}{9+x^2}\,dx$}
\end{frame}

\begin{frame}{Integrales mixtas: identificar el método}
\small Ejercicio 11
\ejercicio{$\displaystyle \int\frac{1}{\sqrt{25-x^2}}\,dx$}
\end{frame}

\begin{frame}{Integrales mixtas: identificar el método}
\small Ejercicio 12
\ejercicio{$\displaystyle \int\left(x^{3/2}+x^{-1/2}\right)dx$}
\end{frame}

\begin{frame}{Integrales mixtas: identificar el método}
\small Ejercicio 13
\ejercicio{$\displaystyle \int 7^{2x}\,dx$}
\end{frame}

\begin{frame}{Integrales mixtas: identificar el método}
\small Ejercicio 14
\ejercicio{$\displaystyle \int\frac{4x^3}{2+x^4}\,dx$}
\end{frame}

\begin{frame}{Integrales mixtas: identificar el método}
\small Ejercicio 15
\ejercicio{$\displaystyle \int\sin x\cos^6x\,dx$}
\end{frame}

\begin{frame}{Integrales mixtas: identificar el método}
\small Ejercicio 16
\ejercicio{$\displaystyle \int\frac{e^x}{4+e^x}\,dx$}
\end{frame}

\section{Retos de cierre}
\begin{frame}{Retos de cierre}
\small Ejercicio 1
\ejercicio{$\displaystyle \int\frac{x^5+2x^3-x}{x^2}\,dx$}
\end{frame}

\begin{frame}{Retos de cierre}
\small Ejercicio 2
\ejercicio{$\displaystyle \int\frac{x^2}{\sqrt{1+x^3}}\,dx$}
\end{frame}

\begin{frame}{Retos de cierre}
\small Ejercicio 3
\ejercicio{$\displaystyle \int\frac{2x-3}{x^2-3x+10}\,dx$}
\end{frame}

\begin{frame}{Retos de cierre}
\small Ejercicio 4
\ejercicio{$\displaystyle \int e^{-3x}(1+e^{-3x})^7\,dx$}
\end{frame}

\begin{frame}{Retos de cierre}
\small Ejercicio 5
\ejercicio{$\displaystyle \int\frac{\cos(2x)}{3+\sin(2x)}\,dx$}
\end{frame}

\begin{frame}{Retos de cierre}
\small Ejercicio 6
\ejercicio{$\displaystyle \int\frac{1}{16+25x^2}\,dx$}
\end{frame}

\begin{frame}{Retos de cierre}
\small Ejercicio 7
\ejercicio{$\displaystyle \int\frac{1}{\sqrt{36-4x^2}}\,dx$}
\end{frame}

\begin{frame}{Retos de cierre}
\small Ejercicio 8
\ejercicio{$\displaystyle \int\left(2\sin^2x+3\cos^2x\right)dx$}
\end{frame}


%=========================================================
% Potencias de seno
%=========================================================
\section{Potencias de seno}
\begin{frame}{Potencias de seno}
\small Ejercicio 1
\ejercicio{$\displaystyle \int \sin^3 x\,dx$}
\end{frame}
\begin{frame}{Potencias de seno}
\small Ejercicio 2
\ejercicio{$\displaystyle \int \sin^5 x\,dx$}
\end{frame}
\begin{frame}{Potencias de seno}
\small Ejercicio 3
\ejercicio{$\displaystyle \int \sin^7 x\,dx$}
\end{frame}
\begin{frame}{Potencias de seno}
\small Ejercicio 4
\ejercicio{$\displaystyle \int \sin^2 x\,dx$}
\end{frame}
\begin{frame}{Potencias de seno}
\small Ejercicio 5
\ejercicio{$\displaystyle \int \sin^4 x\,dx$}
\end{frame}
\begin{frame}{Potencias de seno}
\small Ejercicio 6
\ejercicio{$\displaystyle \int \sin^6 x\,dx$}
\end{frame}
\begin{frame}{Potencias de seno}
\small Ejercicio 7
\ejercicio{$\displaystyle \int \sin^3(2x)\,dx$}
\end{frame}
\begin{frame}{Potencias de seno}
\small Ejercicio 8
\ejercicio{$\displaystyle \int \sin^4(3x)\,dx$}
\end{frame}

%=========================================================
% Potencias de coseno
%=========================================================
\section{Potencias de coseno}
\begin{frame}{Potencias de coseno}
\small Ejercicio 1
\ejercicio{$\displaystyle \int \cos^3 x\,dx$}
\end{frame}
\begin{frame}{Potencias de coseno}
\small Ejercicio 2
\ejercicio{$\displaystyle \int \cos^5 x\,dx$}
\end{frame}
\begin{frame}{Potencias de coseno}
\small Ejercicio 3
\ejercicio{$\displaystyle \int \cos^7 x\,dx$}
\end{frame}
\begin{frame}{Potencias de coseno}
\small Ejercicio 4
\ejercicio{$\displaystyle \int \cos^2 x\,dx$}
\end{frame}
\begin{frame}{Potencias de coseno}
\small Ejercicio 5
\ejercicio{$\displaystyle \int \cos^4 x\,dx$}
\end{frame}
\begin{frame}{Potencias de coseno}
\small Ejercicio 6
\ejercicio{$\displaystyle \int \cos^6 x\,dx$}
\end{frame}
\begin{frame}{Potencias de coseno}
\small Ejercicio 7
\ejercicio{$\displaystyle \int \cos^3(2x)\,dx$}
\end{frame}
\begin{frame}{Potencias de coseno}
\small Ejercicio 8
\ejercicio{$\displaystyle \int \cos^4(3x)\,dx$}
\end{frame}

%=========================================================
% Productos de potencias seno-coseno
%=========================================================
\section{Productos de potencias de seno y coseno}
\begin{frame}{Productos de potencias de seno y coseno}
\small Ejercicio 1
\ejercicio{$\displaystyle \int \sin^3 x\cos^2 x\,dx$}
\end{frame}
\begin{frame}{Productos de potencias de seno y coseno}
\small Ejercicio 2
\ejercicio{$\displaystyle \int \sin^2 x\cos^3 x\,dx$}
\end{frame}
\begin{frame}{Productos de potencias de seno y coseno}
\small Ejercicio 3
\ejercicio{$\displaystyle \int \sin^5 x\cos^4 x\,dx$}
\end{frame}
\begin{frame}{Productos de potencias de seno y coseno}
\small Ejercicio 4
\ejercicio{$\displaystyle \int \sin^4 x\cos^5 x\,dx$}
\end{frame}
\begin{frame}{Productos de potencias de seno y coseno}
\small Ejercicio 5
\ejercicio{$\displaystyle \int \sin^2 x\cos^2 x\,dx$}
\end{frame}
\begin{frame}{Productos de potencias de seno y coseno}
\small Ejercicio 6
\ejercicio{$\displaystyle \int \sin^4 x\cos^2 x\,dx$}
\end{frame}
\begin{frame}{Productos de potencias de seno y coseno}
\small Ejercicio 7
\ejercicio{$\displaystyle \int \sin^2 x\cos^4 x\,dx$}
\end{frame}
\begin{frame}{Productos de potencias de seno y coseno}
\small Ejercicio 8
\ejercicio{$\displaystyle \int \sin^4 x\cos^4 x\,dx$}
\end{frame}

%=========================================================
% Sustitución trigonométrica
%=========================================================
\section{Sustitución trigonométrica}
\begin{frame}{Sustitución trigonométrica}
\small Ejercicio 1
\ejercicio{$\displaystyle \int \frac{dx}{\sqrt{9-x^2}}$}
\end{frame}
\begin{frame}{Sustitución trigonométrica}
\small Ejercicio 2
\ejercicio{$\displaystyle \int \sqrt{16-x^2}\,dx$}
\end{frame}
\begin{frame}{Sustitución trigonométrica}
\small Ejercicio 3
\ejercicio{$\displaystyle \int \frac{x^2}{\sqrt{25-x^2}}\,dx$}
\end{frame}
\begin{frame}{Sustitución trigonométrica}
\small Ejercicio 4
\ejercicio{$\displaystyle \int \frac{dx}{\sqrt{x^2+9}}$}
\end{frame}
\begin{frame}{Sustitución trigonométrica}
\small Ejercicio 5
\ejercicio{$\displaystyle \int \sqrt{x^2+4}\,dx$}
\end{frame}
\begin{frame}{Sustitución trigonométrica}
\small Ejercicio 6
\ejercicio{$\displaystyle \int \frac{x^2}{\sqrt{x^2+16}}\,dx$}
\end{frame}
\begin{frame}{Sustitución trigonométrica}
\small Ejercicio 7
\ejercicio{$\displaystyle \int \frac{dx}{\sqrt{x^2-9}}$}
\end{frame}
\begin{frame}{Sustitución trigonométrica}
\small Ejercicio 8
\ejercicio{$\displaystyle \int \frac{\sqrt{x^2-25}}{x}\,dx$}
\end{frame}

%=========================================================
% Funciones trigonométricas inversas
%=========================================================
\section{Integrales que conducen a $\arcsin$}
\begin{frame}{Integrales que conducen a $\arcsin$}
\small Ejercicio 1
\ejercicio{$\displaystyle \int \frac{dx}{\sqrt{1-x^2}}$}
\end{frame}
\begin{frame}{Integrales que conducen a $\arcsin$}
\small Ejercicio 2
\ejercicio{$\displaystyle \int \frac{dx}{\sqrt{4-x^2}}$}
\end{frame}
\begin{frame}{Integrales que conducen a $\arcsin$}
\small Ejercicio 3
\ejercicio{$\displaystyle \int \frac{dx}{\sqrt{25-4x^2}}$}
\end{frame}
\begin{frame}{Integrales que conducen a $\arcsin$}
\small Ejercicio 4
\ejercicio{$\displaystyle \int \frac{3\,dx}{\sqrt{16-9x^2}}$}
\end{frame}
\section{Integrales que conducen a $\arctan$}
\begin{frame}{Integrales que conducen a $\arctan$}
\small Ejercicio 5
\ejercicio{$\displaystyle \int \frac{dx}{1+x^2}$}
\end{frame}
\begin{frame}{Integrales que conducen a $\arctan$}
\small Ejercicio 6
\ejercicio{$\displaystyle \int \frac{dx}{9+x^2}$}
\end{frame}
\begin{frame}{Integrales que conducen a $\arctan$}
\small Ejercicio 7
\ejercicio{$\displaystyle \int \frac{dx}{4+9x^2}$}
\end{frame}
\begin{frame}{Integrales que conducen a $\arctan$}
\small Ejercicio 8
\ejercicio{$\displaystyle \int \frac{5\,dx}{25+4x^2}$}
\end{frame}

\section{Integración por partes}

\begin{frame}{Integración por partes}
\small Ejercicio 1
\ejercicio{$\displaystyle \int x e^x\,dx$}
\end{frame}

\begin{frame}{Integración por partes}
\small Ejercicio 2
\ejercicio{$\displaystyle \int x\sin x\,dx$}
\end{frame}

\begin{frame}{Integración por partes}
\small Ejercicio 3
\ejercicio{$\displaystyle \int x\cos x\,dx$}
\end{frame}

\begin{frame}{Integración por partes}
\small Ejercicio 4
\ejercicio{$\displaystyle \int x\ln x\,dx$}
\end{frame}

\begin{frame}{Integración por partes}
\small Ejercicio 5
\ejercicio{$\displaystyle \int x^2e^x\,dx$}
\end{frame}

\begin{frame}{Integración por partes}
\small Ejercicio 6
\ejercicio{$\displaystyle \int x^2\sin x\,dx$}
\end{frame}

\begin{frame}{Integración por partes}
\small Ejercicio 7
\ejercicio{$\displaystyle \int \ln x\,dx$}
\end{frame}

\begin{frame}{Integración por partes}
\small Ejercicio 8
\ejercicio{$\displaystyle \int e^x\cos x\,dx$}
\end{frame}

\section{Potencias de tangente y secante}

\begin{frame}{Potencias de tangente y secante}
\small Ejercicio 1
\ejercicio{$\displaystyle \int \tan^3x\,dx$}
\end{frame}

\begin{frame}{Potencias de tangente y secante}
\small Ejercicio 2
\ejercicio{$\displaystyle \int \sec^3x\,dx$}
\end{frame}

\begin{frame}{Potencias de tangente y secante}
\small Ejercicio 3
\ejercicio{$\displaystyle \int \tan^2x\sec^2x\,dx$}
\end{frame}

\begin{frame}{Potencias de tangente y secante}
\small Ejercicio 4
\ejercicio{$\displaystyle \int \tan^4x\sec^2x\,dx$}
\end{frame}

\begin{frame}{Potencias de tangente y secante}
\small Ejercicio 5
\ejercicio{$\displaystyle \int \sec^4x\tan x\,dx$}
\end{frame}

\begin{frame}{Potencias de tangente y secante}
\small Ejercicio 6
\ejercicio{$\displaystyle \int \sec^2x\tan^3x\,dx$}
\end{frame}

\begin{frame}{Potencias de tangente y secante}
\small Ejercicio 7
\ejercicio{$\displaystyle \int \tan^4x\,dx$}
\end{frame}

\begin{frame}{Potencias de tangente y secante}
\small Ejercicio 8
\ejercicio{$\displaystyle \int \sec^5x\,dx$}
\end{frame}

\section{Potencias de cotangente y cosecante}

\begin{frame}{Potencias de cotangente y cosecante}
\small Ejercicio 1
\ejercicio{$\displaystyle \int \cot^3x\,dx$}
\end{frame}

\begin{frame}{Potencias de cotangente y cosecante}
\small Ejercicio 2
\ejercicio{$\displaystyle \int \csc^3x\,dx$}
\end{frame}

\begin{frame}{Potencias de cotangente y cosecante}
\small Ejercicio 3
\ejercicio{$\displaystyle \int \cot^2x\csc^2x\,dx$}
\end{frame}

\begin{frame}{Potencias de cotangente y cosecante}
\small Ejercicio 4
\ejercicio{$\displaystyle \int \cot^4x\csc^2x\,dx$}
\end{frame}

\begin{frame}{Potencias de cotangente y cosecante}
\small Ejercicio 5
\ejercicio{$\displaystyle \int \csc^4x\cot x\,dx$}
\end{frame}

\begin{frame}{Potencias de cotangente y cosecante}
\small Ejercicio 6
\ejercicio{$\displaystyle \int \csc^2x\cot^3x\,dx$}
\end{frame}

\begin{frame}{Potencias de cotangente y cosecante}
\small Ejercicio 7
\ejercicio{$\displaystyle \int \cot^4x\,dx$}
\end{frame}

\begin{frame}{Potencias de cotangente y cosecante}
\small Ejercicio 8
\ejercicio{$\displaystyle \int \csc^5x\,dx$}
\end{frame}

\section{Productos trigonométricos con ángulos distintos}

\begin{frame}{Productos trigonométricos con ángulos distintos}
\small Ejercicio 1
\ejercicio{$\displaystyle \int \sin(3x)\cos(5x)\,dx$}
\end{frame}

\begin{frame}{Productos trigonométricos con ángulos distintos}
\small Ejercicio 2
\ejercicio{$\displaystyle \int \cos(2x)\cos(7x)\,dx$}
\end{frame}

\begin{frame}{Productos trigonométricos con ángulos distintos}
\small Ejercicio 3
\ejercicio{$\displaystyle \int \sin(4x)\sin(2x)\,dx$}
\end{frame}

\begin{frame}{Productos trigonométricos con ángulos distintos}
\small Ejercicio 4
\ejercicio{$\displaystyle \int \sin x\cos(3x)\,dx$}
\end{frame}

\begin{frame}{Productos trigonométricos con ángulos distintos}
\small Ejercicio 5
\ejercicio{$\displaystyle \int \cos(5x)\cos x\,dx$}
\end{frame}

\begin{frame}{Productos trigonométricos con ángulos distintos}
\small Ejercicio 6
\ejercicio{$\displaystyle \int \sin(6x)\sin(3x)\,dx$}
\end{frame}

\begin{frame}{Productos trigonométricos con ángulos distintos}
\small Ejercicio 7
\ejercicio{$\displaystyle \int \sin(2x)\cos(4x)\,dx$}
\end{frame}

\begin{frame}{Productos trigonométricos con ángulos distintos}
\small Ejercicio 8
\ejercicio{$\displaystyle \int \cos(3x)\sin(7x)\,dx$}
\end{frame}

\section{Funciones racionales y división de polinomios}

\begin{frame}{Funciones racionales y división de polinomios}
\small Ejercicio 1
\ejercicio{$\displaystyle \int \frac{x^3+2x^2-x+1}{x^2+1}\,dx$}
\end{frame}

\begin{frame}{Funciones racionales y división de polinomios}
\small Ejercicio 2
\ejercicio{$\displaystyle \int \frac{x^2+3x+5}{x+1}\,dx$}
\end{frame}

\begin{frame}{Funciones racionales y división de polinomios}
\small Ejercicio 3
\ejercicio{$\displaystyle \int \frac{2x^3-x+4}{x^2-1}\,dx$}
\end{frame}

\begin{frame}{Funciones racionales y división de polinomios}
\small Ejercicio 4
\ejercicio{$\displaystyle \int \frac{x^4+1}{x^2+1}\,dx$}
\end{frame}

\begin{frame}{Funciones racionales y división de polinomios}
\small Ejercicio 5
\ejercicio{$\displaystyle \int \frac{3x^3+2x^2+1}{x^2+x}\,dx$}
\end{frame}

\begin{frame}{Funciones racionales y división de polinomios}
\small Ejercicio 6
\ejercicio{$\displaystyle \int \frac{x^3-1}{x-1}\,dx$}
\end{frame}

\begin{frame}{Funciones racionales y división de polinomios}
\small Ejercicio 7
\ejercicio{$\displaystyle \int \frac{x^4-x^2+1}{x^2+1}\,dx$}
\end{frame}

\begin{frame}{Funciones racionales y división de polinomios}
\small Ejercicio 8
\ejercicio{$\displaystyle \int \frac{2x^3+5x^2-3}{x+2}\,dx$}
\end{frame}

\section{Fracciones parciales}

\begin{frame}{Fracciones parciales}
\small Ejercicio 1
\ejercicio{$\displaystyle \int \frac{dx}{(x-1)(x+2)}$}
\end{frame}

\begin{frame}{Fracciones parciales}
\small Ejercicio 2
\ejercicio{$\displaystyle \int \frac{3x+5}{(x-1)(x+3)}\,dx$}
\end{frame}

\begin{frame}{Fracciones parciales}
\small Ejercicio 3
\ejercicio{$\displaystyle \int \frac{dx}{x(x+1)^2}$}
\end{frame}

\begin{frame}{Fracciones parciales}
\small Ejercicio 4
\ejercicio{$\displaystyle \int \frac{2x+1}{(x-2)^2(x+1)}\,dx$}
\end{frame}

\begin{frame}{Fracciones parciales}
\small Ejercicio 5
\ejercicio{$\displaystyle \int \frac{dx}{(x^2+1)(x-1)}$}
\end{frame}

\begin{frame}{Fracciones parciales}
\small Ejercicio 6
\ejercicio{$\displaystyle \int \frac{x^2+1}{x(x-1)(x+2)}\,dx$}
\end{frame}

\begin{frame}{Fracciones parciales}
\small Ejercicio 7
\ejercicio{$\displaystyle \int \frac{dx}{(x^2+4)(x^2+1)}$}
\end{frame}

\begin{frame}{Fracciones parciales}
\small Ejercicio 8
\ejercicio{$\displaystyle \int \frac{x+1}{(x^2+1)^2}\,dx$}
\end{frame}

\section{Combinación de técnicas}

\begin{frame}{Combinación de técnicas}
\small Ejercicio 1
\ejercicio{$\displaystyle \int x\arctan x\,dx$}
\end{frame}

\begin{frame}{Combinación de técnicas}
\small Ejercicio 2
\ejercicio{$\displaystyle \int x^2\ln x\,dx$}
\end{frame}

\begin{frame}{Combinación de técnicas}
\small Ejercicio 3
\ejercicio{$\displaystyle \int x e^{x^2}\cos(e^{x^2})\,dx$}
\end{frame}

\begin{frame}{Combinación de técnicas}
\small Ejercicio 4
\ejercicio{$\displaystyle \int \frac{x^3}{\sqrt{9-x^2}}\,dx$}
\end{frame}

\begin{frame}{Combinación de técnicas}
\small Ejercicio 5
\ejercicio{$\displaystyle \int \frac{x^2}{x^2+4}\,dx$}
\end{frame}

\begin{frame}{Combinación de técnicas}
\small Ejercicio 6
\ejercicio{$\displaystyle \int e^x\sin x\,dx$}
\end{frame}

\begin{frame}{Combinación de técnicas}
\small Ejercicio 7
\ejercicio{$\displaystyle \int \frac{x^2+1}{x^3+x}\,dx$}
\end{frame}

\begin{frame}{Combinación de técnicas}
\small Ejercicio 8
\ejercicio{$\displaystyle \int \sin^3x\cos^4x\,dx$}
\end{frame}

\end{document}
